% -------------------------------------------------------------------------
% WBS ID: RES-DAT-GEO
% Deliverable: Bathymetry, Topography and Coastal Geometry
% Status: In progress
% Evidence status: Regional model data requirements established; dataset selection pending
% Review status: Not reviewed
% Last updated: 2026-07-19
% Dependencies: RES-MOD-SPEC, RES-NUM-SPEC, RES-PHY-NRS
% Downstream interfaces: QGIS/geospatial domain lock, RES-VER-SPEC, Software data ingestion
% -------------------------------------------------------------------------

\section{RES-DAT-GEO --- Bathymetry, Topography and Coastal Geometry}
\label{sec:res-dat-geo}

\subsection{Purpose and Scope}
\label{sec:res-dat-geo-purpose}

This Deliverable defines the bathymetry, topography and coastal-geometry
data required to build the regional 2D NLSWE domain, the two
regional-to-local corridor extracts and the local 3D URANS--VOF
terrain surfaces. It covers dataset selection, fallback hierarchy,
coordinate reference systems, vertical datums, preprocessing and
handoff requirements. It does not define arbitrary corridor sketches:
the corridor footprint is a data/model interface constructed from the
event source, target coast, observation footprint, defence geometry and
numerical-boundary requirements.

\subsection{Research Questions}
\label{sec:res-dat-geo-research-questions}

The data decision asks:

\begin{enumerate}
  \item Which topobathymetric source is authoritative inside the Japan
    nearshore and local corridors?
  \item Which global products are acceptable for the outer regional
    domain or fallback coverage?
  \item How are horizontal CRS, vertical datum, sign convention,
    resolution and licence metadata preserved through clipping and
    interpolation?
  \item Which preprocessing uncertainty must be handed to `RES-VER`
    and Software before validation claims are made?
\end{enumerate}

\subsection{Terminology and Definitions}
\label{sec:res-dat-geo-definitions}

% TODO: Define the technical terminology, symbols, quantities and
% classifications used in this Deliverable.

\subsection{Literature Evidence}
\label{sec:res-dat-geo-literature-evidence}

% TODO: Record evidence from peer-reviewed papers, authoritative datasets,
% standards, technical reports and primary documentation.
%
% Do not create a detached literature-summary list. Explain how each source
% supports, contradicts or limits a project decision.

\textcite{jodcJegg500} documents J-EGG500 as a 500 m gridded
bathymetric dataset around Japan compiled from depth measurements from
Japanese marine survey institutions. Source role: supports the primary
Japan-region bathymetry candidate. Project-specific conclusion:
J-EGG500 is the preferred first candidate for offshore bathymetry in
the Kamaishi and Sendai corridors if licence, datum and import metadata
are acceptable. Limitation: the source page notes heterogeneous data
quality and smoothing, so local relief and harbour/breakwater details
shall not be assumed resolved.

\textcite{gebco2026grid} provides the current global 15 arc-second
terrain grid for ocean and land and a Type Identifier grid. Source
role: supports outer-domain and fallback topobathymetry. Project-
specific conclusion: GEBCO is selected as the first global fallback and
outer-domain source, not as the preferred source for local defence
geometry. Limitation: GEBCO's vertical datum is treated as mean sea
level while shallow-water source datums can vary, so nearshore datum
reconciliation remains mandatory.

\textcite{noaa2022etopo} provides NOAA ETOPO 2022 as a 15 arc-second
global relief model designed for applications including tsunami
forecasting and modelling. Source role: supports secondary global
fallback and cross-check. Project-specific conclusion: ETOPO shall be
used to fill or audit outer-domain coverage if J-EGG500/GEBCO coverage
or provenance is insufficient. Limitation: it is not selected for
local barrier-scale topography.

\textcite{gsiDemDownload} documents GSI digital elevation models with
1 m, 5 m and 10 m mesh products. Source role: supports Japan onshore
topography. Project-specific conclusion: GSI DEM is the preferred
onshore topographic source for corridor clipping and inundation/run-up
comparison where available. Limitation: mixed DEM classes have
different survey methods and accuracy; class boundaries shall be
retained in the provenance manifest.

\subsection{Governing Relations and Quantitative Definitions}
\label{sec:res-dat-geo-governing-relations}

% TODO: Record relevant equations, dimensionless groups, empirical models,
% extraction rules or quantitative definitions.
%
% Use align environments for displayed equations.
% Every displayed equation must have a unique \label{eq:...}.
% Do not place punctuation after displayed equations.

\subsection{Required Data Variables and Units}
\label{sec:res-dat-geo-required-data-variables-units}

% TODO: Complete this RES-DAT Domain-specific subsection.

\subsection{Candidate Data Sources}
\label{sec:res-dat-geo-candidate-data-sources}

The selected/fallback hierarchy for the Tōhoku case is:

\begin{enumerate}
  \item \textbf{Japan offshore bathymetry:} J-EGG500
    \parencite{jodcJegg500}, subject to licence, datum and coverage
    checks for the selected corridors.
  \item \textbf{Japan onshore topography:} GSI DEM
    \parencite{gsiDemDownload}, selecting the finest available mesh
    class that is legally usable and spatially continuous over each
    corridor.
  \item \textbf{Outer regional and fallback topobathymetry:}
    GEBCO\_2026 \parencite{gebco2026grid}.
  \item \textbf{Secondary global fallback or cross-check:}
    NOAA ETOPO 2022 \parencite{noaa2022etopo}.
\end{enumerate}

The selected approach is a stitched and provenance-preserving terrain
model rather than one unqualified raster. Each cell or clipped raster
tile shall retain the source product, source resolution, interpolation
method, horizontal CRS, vertical datum and sign convention.

\subsection{Spatial and Temporal Coverage}
\label{sec:res-dat-geo-spatial-temporal-coverage}

Coverage shall include:

\begin{itemize}
  \item the outer Tōhoku regional propagation domain, including
    artificial-boundary and sponge regions;
  \item the epicentre-to-Kamaishi corridor and enough lateral margin to
    contain the severe-impact observations, bay/coastline geometry and
    relaxation zones;
  \item the epicentre-to-Sendai coastal-plain corridor and enough
    inland coverage to include observed inundation limits;
  \item local 3D terrain patches around wall-type and
    obstacle/dissipating barrier classes.
\end{itemize}

Topobathymetric products are treated as static during one hydrodynamic
run. The earthquake source updates the bed through `RES-DAT-SRC`; it
does not change the static terrain-data selection.

\subsection{Coordinate and Datum Requirements}
\label{sec:res-dat-geo-coordinate-datum-requirements}

The regional model uses a horizontal projected coordinate frame
$(x,y)$ and a positive-upward bed elevation $b$. Source datasets shall
record horizontal CRS, vertical datum, units, coastline convention and
any tidal or geoid correction used before interpolation to the
finite-volume mesh.

Base acquisition coordinates shall remain longitude/latitude
$(\lambda,\phi)$ until the corridor target points are finalised in GIS.
The solver-facing terrain shall be delivered in a projected metric CRS
chosen for the final regional and local domains. Any dataset whose
horizontal CRS or vertical datum cannot be confirmed shall be marked
``usable only for visual screening'' until metadata are resolved.

\subsection{Data Processing and Transformation}
\label{sec:res-dat-geo-data-processing-transformation}

The required transformation pipeline is:

\begin{align}
  (\lambda,\phi,z)
  &\rightarrow
  (E,N,b_0)
  \rightarrow
  (x,y,b_0)
  \rightarrow
  b_0^{\mathrm{mesh}}
  \rightarrow
  b_{0,i}
  \rightarrow
  b_i(t)
  \label{eq:res-dat-geo-regional-bed-data-pipeline}
\end{align}

For positive-upward elevation data:

\begin{align}
  b_0
  &=
  z
  -
  z_{\mathrm{ref}}
  \label{eq:res-dat-geo-positive-upward-bed}
\end{align}

For raw depth $d$ positive downward:

\begin{align}
  b_0
  &=
  -d
  -
  z_{\mathrm{ref}}
  \label{eq:res-dat-geo-positive-downward-depth-conversion}
\end{align}

The cell-average static bed is:

\begin{align}
  b_{0,i}
  &=
  \frac{1}{A_i}
  \int_{K_i}
  b_0^{\mathrm{mesh}}(x,y)
  \,\dd A
  \label{eq:res-dat-geo-cell-average-static-bed}
\end{align}

QGIS/geospatial work owns dataset selection, domain placement,
coordinate transformation and source-data preparation. It is a
dependency of the regional model, not part of the mathematical model
itself.

\subsection{Uncertainty, Provenance and Licensing}
\label{sec:res-dat-geo-uncertainty-provenance-licensing}

The terrain manifest shall store, for every source tile or derived
terrain surface:

\begin{itemize}
  \item source product, provider, citation key and access URL;
  \item licence or usage note;
  \item horizontal CRS, vertical datum and sign convention;
  \item native resolution and interpolation or smoothing operations;
  \item clipping polygon, corridor identifier and mesh-projection
    method;
  \item uncertainty flags for datum conversion, coastline alignment,
    resampling, smoothing and source-age heterogeneity.
\end{itemize}

Project conclusion: a terrain raster without this metadata is not an
accepted validation input. Limitation: the metadata requirements do not
quantify the full bathymetric error; they make the uncertainty
auditable by `RES-VER`.

\subsection{Data Acceptance Criteria}
\label{sec:res-dat-geo-data-acceptance-criteria}

% TODO: Complete this RES-DAT Domain-specific subsection.

\subsection{Candidate Approaches}
\label{sec:res-dat-geo-candidate-approaches}

% TODO: Define the credible candidate approaches considered.

\subsection{Critical Comparison}
\label{sec:res-dat-geo-critical-comparison}

% TODO: Compare candidates using physically and project-relevant criteria.
% Do not select an approach solely on popularity or implementation ease.

\subsection{Assumptions and Applicability}
\label{sec:res-dat-geo-assumptions}

% TODO: State assumptions and the conditions under which they are valid.

\subsection{Limitations and Failure Conditions}
\label{sec:res-dat-geo-limitations}

% TODO: State where the evidence, model or selected approach ceases to be
% reliable or applicable.

\subsection{Project-Specific Conclusions}
\label{sec:res-dat-geo-conclusions}

The Tōhoku case-study terrain workflow shall use Japan-specific sources
where available and global products only where they are physically and
numerically appropriate. J-EGG500 and GSI DEM are selected as the
first candidates for Japan offshore and onshore data respectively.
GEBCO\_2026 is selected for outer-domain/fallback coverage, with ETOPO
2022 retained as a secondary fallback and consistency check. The local
barrier mesh shall not infer defence geometry from 15 arc-second global
products.

\subsection{Selected Approach and Rationale}
\label{sec:res-dat-geo-selected-approach}

Use a provenance-preserving terrain hierarchy:

\begin{quote}
  J-EGG500 offshore bathymetry $+$
  GSI onshore DEM where available, merged with GEBCO\_2026 for
  outer-domain and fallback coverage, and ETOPO 2022 as a secondary
  fallback/cross-check
\end{quote}

Source role: supports terrain-data finalisation for the regional model,
the two corridor extracts and local 3D terrain. Supporting citations:
\textcite{jodcJegg500}, \textcite{gsiDemDownload},
\textcite{gebco2026grid} and \textcite{noaa2022etopo}. Downstream
handoff: Software shall ingest these sources through a dataset
manifest, perform coordinate transformation, clip corridor rasters,
interpolate to mesh/field containers and write structured outputs with
source provenance.

\subsection{Unresolved Decisions}
\label{sec:res-dat-geo-unresolved-decisions}

% TODO: Record unresolved questions, required evidence and decision owners.

The mathematical data requirements are established. The actual QGIS
domain lock, final bathymetry/topography dataset, datum reconciliation,
coastline treatment, smoothing policy and mesh-projection workflow
remain open.

Open source-specific decisions are the final J-EGG500 and GSI licence
acceptance, any required Japanese geodetic/vertical datum conversion,
GEBCO/ETOPO blending boundary, and whether higher-resolution
site-specific harbour or defence surveys can be obtained for Kamaishi.
TODO--missing-source: no dedicated local Kamaishi bay-mouth breakwater
geometry or survey source has yet been confirmed in the WBS
bibliography.

\subsection{Requirements and Handoffs}
\label{sec:res-dat-geo-requirements}

% TODO: State requirements passed to other Research Domains, Software,
% verification activities or later execution work.

The data workstream shall provide mesh-projected static bed values,
source-bed increments where applicable, domain boundaries, sponge-layer
candidate regions and provenance metadata before validation or
implementation can be declared complete.

Software-facing requirements now begin in parallel with remaining
research: dataset manifest, coordinate-transform workflow,
corridor-polygon generation, raster/vector clipping,
bathymetry/topography interpolation, HDF5 or equivalent structured
terrain output, and replay-boundary metadata for the local 3D inlet.

\subsection{Deliverable Completion Record}
\label{sec:res-dat-geo-completion-record}

% TODO: Record whether the Deliverable is:
%
% - Not started
% - In progress
% - Draft complete
% - Reviewed
% - Accepted
%
% Acceptance requires:
%
% - the research questions to be answered;
% - claims to be supported by appropriate evidence;
% - assumptions and limitations to be explicit;
% - the project-specific conclusion to be stated;
% - downstream requirements to be traceable;
% - unresolved decisions to be closed or formally deferred.
